\section{Mathematical Framework (Experimental Setup)}\label{sec:c7s2}
\lipsum[1-1]

This section builds on the material in \cref{sec:c7s1}, \cref{ch:8} and the prior work of \citet{main-36}; see also \citep{main-15,main-16}.

\begin{align}\label{eq:c7s2-a}
  \mathbf{x}_{k+1} &= \mathbf{A}\mathbf{x}_k + \mathbf{B}\mathbf{u}_k + \mathbf{w}_k, & \mathbf{w}_k&\sim\mathcal{N}(0,\mathbf{Q}),\\
  \mathbf{y}_k &= \mathbf{C}\mathbf{x}_k + \mathbf{v}_k, & \mathbf{v}_k&\sim\mathcal{N}(0,\mathbf{R}).\label{eq:c7s2-b}
\end{align}

Equations~\eqref{eq:c7s2-a} and~\eqref{eq:c7s2-b} are used throughout the remainder of this section.

\kant[1-10]

\subsection{Quantitative behaviour}\label{sec:c7s2-q}
Results are collected in \cref{tab:c7s2} and visualised in \cref{fig:c7s2-f1}. Similar trends were reported in~\citep{main-27,main-10}.

\begin{table}[htbp]
\centering
\caption{Summary of measured quantities for experiment set \texttt{c7s2}.}\label{tab:c7s2}
\begin{tabular}{lrrr}
\toprule
Method & Score & Latency & Error \\
\midrule
Alpha & \num{28.5} & \SI{6.15}{\milli\second} & \num{7.797e-5} \\
Beta & \num{96.2} & \SI{8.17}{\milli\second} & \num{1.926e-5} \\
Gamma & \num{18.2} & \SI{3.31}{\milli\second} & \num{4.068e-2} \\
Delta & \num{36.7} & \SI{8.45}{\milli\second} & \num{2.231e-4} \\
Epsilon & \num{66.4} & \SI{8.01}{\milli\second} & \num{8.637e-5} \\
\bottomrule
\end{tabular}
\end{table}

\begin{figure}[htbp]
\centering
\begin{tikzpicture}\begin{groupplot}[group style={group size=2 by 1,horizontal sep=1.8cm},width=0.43\linewidth,height=5cm]
\nextgroupplot[title={(a)},domain=-3:3]\addplot[teal,thick,samples=200]{x^3-4*x};
\nextgroupplot[title={(b)},domain=-3:3]\addplot[black,thick,samples=200]{4*exp(-x^2)*cos(deg(4*x))};
\end{groupplot}\end{tikzpicture}
\caption{Generated figure for \texttt{c7s2-f1} (parameters $a=4$, $b=4$).}\label{fig:c7s2-f1}
\end{figure}

\lipsum[48-49]

\subsection{Implementation notes}\label{sec:c7s2-i}
The core routine is shown in \cref{lst:c7s2}; its complexity is $\mathcal{O}(N\log N)$ per iteration~\cite{main-25}.

\begin{lstlisting}[language=python,caption={Reference implementation for c7s2.},label={lst:c7s2}]
def conjugate_gradient(A, b, tol=1e-8, maxit=500):
    x = b * 0
    r = b - A @ x
    p = r.copy()
    for k in range(maxit):
        Ap = A @ p
        alpha = (r @ r) / (p @ Ap)
        x += alpha * p
        r_new = r - alpha * Ap
        if (r_new @ r_new) ** 0.5 < tol:
            break
        p = r_new + ((r_new @ r_new) / (r @ r)) * p
        r = r_new
    return x
\end{lstlisting}

\kant[1-11]
\begin{theorem}\label{thm:c7s2}
Let $A\in\mathbb{R}^{n\times n}$ be symmetric positive definite with condition number $\kappa$. Then the iteration converges and
$\lVert e_k\rVert_A \le 2\bigl(\tfrac{\sqrt\kappa-1}{\sqrt\kappa+1}\bigr)^k\lVert e_0\rVert_A$.
\end{theorem}
\begin{enumerate}[label=(\roman*)]
\item the bound in \cref{thm:c7s2} is sharp for some initial errors;
\item the constant $2$ cannot be improved in general;
\item preconditioning reduces the effective $\kappa$.
\end{enumerate}

\kant[4-12]
